\section{Current Measured Evidence}
\label{sec:results}

\subsection{Evidence boundary}

This section contains every numerical result used to characterize v0.1.0.
There is no hidden large-model result and no accelerator benchmark. The
measurements fall into three categories: invariant tests, a small learnability
diagnostic, and a saturated-state CPU timing run.

\subsection{Automated release gate}

The recorded release gate ran on Windows, Python 3.11.15, PyTorch
2.13.0+cpu. It completed bytecode compilation, 50 tests, two training smoke
runs, and one performance smoke run with zero non-zero return codes.

\begin{table}[h]
\centering
\caption{Recorded release-gate outcomes. Smoke runs check executable paths;
their accuracy is not a trained-model result.}
\label{tab:release-gate}
\begin{tabular}{@{}lrrl@{}}
\toprule
Gate & Elapsed (s) & Return code & Interpretation \\
\midrule
Compile source, scripts, tests & 0.152 & 0 & Syntax/import compilation \\
Pytest suite & 5.815 & 0 & 50 tests passed \\
Diagnostic MQAR smoke & 3.368 & 0 & Two optimizer steps completed \\
Zoology-aligned MQAR smoke & 3.164 & 0 & One tiny epoch completed \\
Stream benchmark smoke & 1.971 & 0 & Timing/state path completed \\
\bottomrule
\end{tabular}
\end{table}

The two-step and one-epoch smoke tasks report zero evaluation accuracy on
their tiny held-out samples. That is expected to be possible after almost no
optimization and is preserved in the JSON rather than filtered out. The
release criterion for those jobs is numerical and operational completion, not
task mastery.

\subsection{Small MQAR learnability diagnostic}

A separate diagnostic trained a 14,288-parameter, two-layer model for 300
steps on the ALUCLU-specific next-token MQAR generator. The setup used
\(d=16\), batch size 32, two key-value pairs, two queries, four key classes,
four value classes, learning rate \(3\times10^{-3}\), and seed zero. Sixteen
evaluation batches produced:

\begin{align}
  \text{answer accuracy} &= 0.904296875,\\
  \text{evaluation loss} &= 0.1969321771,\\
  \text{elapsed time} &= 29.9342\ \mathrm{s},\\
  \text{training throughput} &= 2886.33\ \mathrm{tokens/s},\\
  \text{batch-1 terminal state} &= 9200\ \mathrm{bytes}.
\end{align}

This is evidence that the integrated model can learn one deliberately easy
associative-recall distribution. It is not an estimate of long-context
capacity: the vocabularies are tiny, only one seed was run, the relevant
sequence lengths are short, and no baseline comparison is attached.

\subsection{Saturated-state CPU reference timing}

The reference timing configuration used \(d=32\), four layers, vocabulary
8192, 363,385 parameters, one PyTorch CPU thread, 128-token prefill, 560
warm-up tokens, and three repeated 32-token decode measurements. The 560-token
warm-up exceeds the reported 528-token state-saturation requirement, so decode
started with every configured episodic capacity allocated.

\begin{table}[h]
\centering
\caption{Saturated reference-backend timing. Random token generation is
outside timed regions. Values are descriptive for this machine and code path,
not baseline comparisons.}
\label{tab:cpu-results}
\begin{tabular}{@{}lrrr@{}}
\toprule
Metric & Minimum & Median & Maximum \\
\midrule
Decode throughput (tokens/s) & 127.96 & 143.71 & 180.28 \\
Decode latency (ms/token) & -- & 6.9586 & -- \\
Prefill throughput (tokens/s) & 500.30 & 587.36 & 916.27 \\
\bottomrule
\end{tabular}
\end{table}

The saturated logical state payload for batch one was 164,576 bytes. The wide
range across only three samples illustrates why this result cannot support a
systems-leadership claim. The benchmark is nevertheless useful as a
regression anchor because it records thread count, environment, model
configuration, timed-region policy, samples, and saturation status.

\subsection{What can be concluded}

The combined evidence supports four statements.

\begin{enumerate}
  \item The released code paths compile and execute in the recorded CPU
  environment.
  \item The tested causal, state, compression, routing, and equivalence
  invariants hold for the test domain.
  \item The full four-lane model is trainable on a small diagnostic and can
  exceed chance substantially under that protocol.
  \item A fully allocated recurrent state can decode indefinitely without
  persistent tensor payload growing with the number of processed tokens.
\end{enumerate}

It does not yet answer whether the lane allocation is Pareto efficient, whether
the surprise policy predicts future utility, whether the archive improves
language modeling, whether training remains stable at scale, or whether a
specialized implementation outperforms optimized attention or state-space
baselines.
